% pauli_exclusion_ML_slides.tex
% Beamer presentation: "Pauli Was Doing Machine Learning"
% Thomas J. Sargent, 2026
% ~30 minutes, ~24 slides

\documentclass[aspectratio=169,12pt]{beamer}

\usetheme{Madrid}
\usecolortheme{seahorse}
\setbeamertemplate{navigation symbols}{}
\setbeamertemplate{footline}[frame number]

\usepackage{amsmath,amssymb}
\usepackage{booktabs}
\usepackage{array}
\usepackage{multirow}
\usepackage{microtype}
\usepackage{hyperref}
\usepackage{xcolor}

\hypersetup{colorlinks=true, linkcolor=blue, citecolor=blue, urlcolor=blue}

\definecolor{theoryblue}{RGB}{0,80,160}
\definecolor{datared}{RGB}{180,30,30}
\definecolor{mlgreen}{RGB}{30,130,60}

\title[Pauli Was Doing Machine Learning]{%
  \textbf{Pauli Was Doing Machine Learning}\\[4pt]
  {\normalsize The Exclusion Principle as a Parsimonious\\
   Four-Parameter Model}}
\author{Thomas J.\ Sargent}
\institute{New York University \& Hoover Institution}
\date{2026}

%----------------------------------------------------------------------
\begin{document}
%----------------------------------------------------------------------

\begin{frame}[plain]
  \titlepage
\end{frame}

%----------------------------------------------------------------------
\begin{frame}{The Central Claim}

\begin{block}{Thesis}
When Pauli proposed the exclusion principle in 1925 to organise the
complex structure of atomic spectra, he was---in the language we now
use---doing \textbf{machine learning} and \textbf{statistical inference}.
\end{block}

\bigskip
\textbf{Three acts:}
\begin{enumerate}
  \item[\color{theoryblue}1.] \textbf{Empirical pattern} (1913--1924):
    experiments reveal that shell capacities, spectral multiplicities, and
    Zeeman splittings cannot be explained by three quantum numbers
  \item[\color{datared}2.] \textbf{Stoner's consolidation} (1924):
    X-ray data show every magnetic sublevel holds exactly \emph{two} electrons
  \item[\color{mlgreen}3.] \textbf{Pauli's minimal model} (1925):
    a fourth, two-valued quantum number plus one binary constraint predicts
    \emph{all} observed electron capacities with \emph{zero} free parameters
\end{enumerate}

\bigskip
\small The episode illuminates---even more starkly than the photoelectric
effect---the power of a single theoretical postulate to compress an
enormous combinatorial problem into a counting exercise.

\end{frame}

%----------------------------------------------------------------------
\begin{frame}{Background: Building Blocks of Atomic Structure, 1913--1924}

\begin{columns}[T]
\begin{column}{0.55\textwidth}
\textbf{What was known (Bohr 1913; Sommerfeld 1916):}
\begin{itemize}
  \item \textbf{Principal} quantum number~$n$: shell energy
    $E_n = -13.6\,\text{eV}/n^2$
  \item \textbf{Azimuthal} quantum number~$l\in\{0,\ldots,n{-}1\}$:
    orbital angular momentum
  \item \textbf{Magnetic} quantum number~$m_l\in\{-l,\ldots,+l\}$:
    projection on field axis
\end{itemize}

\bigskip
\textbf{States per shell} from three numbers:
\[
  \sum_{l=0}^{n-1}(2l+1) = n^2
\]
\end{column}
\begin{column}{0.42\textwidth}
\begin{alertblock}{The Core Puzzle}
\begin{tabular}{cccc}
\toprule
Shell & $n$ & Predicted & Observed \\
      &     & $n^2$     & $2n^2$ \\
\midrule
K & 1 &  1 &  2 \\
L & 2 &  4 &  8 \\
M & 3 &  9 & 18 \\
N & 4 & 16 & 32 \\
\bottomrule
\end{tabular}

\medskip
Observed capacity is \textbf{exactly twice}
the predicted capacity.\\[4pt]
\emph{Why?}
\end{alertblock}
\end{column}
\end{columns}

\end{frame}

%----------------------------------------------------------------------
\begin{frame}[shrink=15]{Three Failures of the Three-Quantum-Number Scheme}

\begin{columns}[T]
\begin{column}{0.52\textwidth}
\textbf{1. Shell closures unexplained}

\smallskip
Bohr proposed noble gases = closed shells:
He (2), Ne (10), Ar (18), Kr (36).
\emph{No dynamical principle explains why
a shell closes at $2n^2$, not $n^2$.}

\bigskip
\textbf{2. Anomalous Zeeman effect}

\smallskip
Most spectral lines split into \textbf{more
than 3 components} in a magnetic field.
Classical theory (Lorentz) predicts exactly~3.
Pauli (1924): ``\emph{How can one look happy when
he is thinking about the anomalous Zeeman
effect?}''

\bigskip
\textbf{3. Half-integer angular momenta}

\smallskip
Landé's (1921) empirical g-factor formula
requires total angular momentum $J$ to take
\textbf{half-integer} values---impossible with
three integer quantum numbers.
\end{column}
\begin{column}{0.45\textwidth}
\begin{block}{Landé's Empirical Formula (1921)}
\[
  g_J = 1 + \frac{J(J+1)+S(S+1)-L(L+1)}{2J(J+1)}
\]
Fits all observed Zeeman spacings perfectly.
\smallskip

Observed g-factors for sodium:
\begin{align*}
  g({}^2S_{1/2}) &= 2.00 \\
  g({}^2P_{1/2}) &= 0.67 \\
  g({}^2P_{3/2}) &= 1.33
\end{align*}
Non-integer values; require $g_s = 2$.
\end{block}
\end{column}
\end{columns}

\end{frame}

%----------------------------------------------------------------------
\begin{frame}[shrink=15]{Stoner's 1924 Key Finding: The Decisive Empirical Clue}

\begin{columns}[T]
\begin{column}{0.55\textwidth}
Edmund Stoner used X-ray spectroscopic level data to count electrons
per \emph{sublevel} of each shell.

\bigskip
\textbf{Stoner's observation:} in a magnetic field, the $m_l$ sublevels
separate.  X-ray data show each $m_l$ sublevel holds \textbf{two} electrons.

\bigskip
\[
  \text{capacity}(n,l) = 2\times(2l+1)
\]

Summing: capacity$(n) = 2n^2$.~\checkmark

\bigskip
\textbf{But: why exactly two per $m_l$ sublevel?}

Stoner could not answer this.
\end{column}
\begin{column}{0.42\textwidth}
\begin{alertblock}{Stoner (1924) as Precursor Data}
\begin{tabular}{ccc}
\toprule
Subshell & $2(2l+1)$ & Observed \\
\midrule
$s\;(l=0)$ & $2$ & $2$ \\
$p\;(l=1)$ & $6$ & $6$ \\
$d\;(l=2)$ & $10$ & $10$ \\
$f\;(l=3)$ & $14$ & $14$ \\
\bottomrule
\end{tabular}
\end{alertblock}

\bigskip
\begin{block}{ML Language}
Stoner performed the analogue of Lenard's
feature-importance experiment: he showed
that the missing factor is \textbf{exactly~2}
per sublevel---a binary quantity.
\end{block}
\end{column}
\end{columns}

\end{frame}

%----------------------------------------------------------------------
\begin{frame}[shrink=40]{Pauli's 1925 Insight: Model Architecture}

\begin{columns}[T]
\begin{column}{0.52\textwidth}
\textbf{Pauli's postulate (January 1925):}

\bigskip
Augment the state descriptor with a
\textbf{fourth quantum number}~$m_s$:
\[
  m_s \in \bigl\{+\tfrac{1}{2},\;-\tfrac{1}{2}\bigr\}
\]
Pauli called this a \emph{Zweideutigkeit}---a
classically non-describable two-valuedness.

\bigskip
Then assert the \textbf{exclusion constraint}:
\[
  \boxed{n_s \in \{0,\,1\}}
\]
No two electrons share the same 4-tuple
$(n,\,l,\,m_l,\,m_s)$.
\end{column}
\begin{column}{0.45\textwidth}
\begin{alertblock}{The Prediction (1925)}
Each subshell $(n,l)$ has
\[
  \underbrace{(2l+1)}_{m_l\text{ values}}
  \times
  \underbrace{2}_{m_s\text{ values}}
  = 2(2l+1)
\]
maximum electrons.

\medskip
Shell capacity:
\[
  \sum_{l=0}^{n-1}2(2l+1) = 2n^2
\]

\medskip
\textbf{Written before a physical
interpretation of $m_s$ existed!}
\end{alertblock}
\end{column}
\end{columns}

\bigskip
\begin{block}{The Inductive Bias}
The four quantum numbers + exclusion constraint restrict the model class from
\emph{all possible occupation rules} to a \textbf{unique counting function} with
zero adjustable parameters.
\end{block}

\end{frame}

%----------------------------------------------------------------------
\begin{frame}[shrink=20]{The Four-Parameter State Space}

\textbf{Each electron's state is a 4-tuple $(n,\,l,\,m_l,\,m_s)$:}

\medskip
\begin{center}
\renewcommand{\arraystretch}{1.2}
\begin{tabular}{clp{4.5cm}p{4.5cm}}
\toprule
\# & Quantum number & Range & Identifying observation \\
\midrule
1 & $n$ (principal) & $1,2,3,\ldots$ & Chemical periodicity; discrete X-ray edges \\
2 & $l$ (azimuthal) & $0,1,\ldots,n-1$ & Fine structure of X-ray levels; Sommerfeld \\
3 & $m_l$ (magnetic) & $-l,\ldots,+l$ & Normal Zeeman effect (3 lines) \\
4 & $m_s$ (spin) & $+\tfrac{1}{2},-\tfrac{1}{2}$ & Stoner's factor of~2; alkali doublets \\
\midrule
\multicolumn{2}{l}{Exclusion: $n_s \leq 1$} & & Noble-gas stability; closed shells \\
\bottomrule
\end{tabular}
\end{center}

\bigskip
\begin{columns}[T]
\begin{column}{0.48\textwidth}
\begin{block}{States per shell (three quantum numbers)}
$\displaystyle\sum_{l=0}^{n-1}(2l+1) = n^2$
\quad\textcolor{datared}{--- half of observed}
\end{block}
\end{column}
\begin{column}{0.48\textwidth}
\begin{alertblock}{States per shell (four quantum numbers)}
$\displaystyle\sum_{l=0}^{n-1}2(2l+1) = 2n^2$
\quad\textcolor{mlgreen}{--- exactly observed}
\end{alertblock}
\end{column}
\end{columns}

\end{frame}

%----------------------------------------------------------------------
\begin{frame}[shrink=15]{The Empirical Dataset: Three Classes of Observations}

\begin{columns}[T]
\begin{column}{0.48\textwidth}
\textbf{Class~1: Subshell electron capacities}
(X-ray spectroscopy + noble-gas chemistry)

\medskip
\begin{tabular}{ccr}
\toprule
Subshell & $l$ & Observed $e^-$ \\
\midrule
$1s$ & 0 & 2 \\
$2s$ & 0 & 2 \\
$2p$ & 1 & 6 \\
$3s$ & 0 & 2 \\
$3p$ & 1 & 6 \\
$3d$ & 2 & 10 \\
$4s$--$4f$ & 0--3 & 2,6,10,14 \\
\bottomrule
\end{tabular}

\bigskip
\textbf{Class~2: Spectral term multiplicities}
\begin{itemize}
  \item Alkali metals: doublets ($S=\tfrac{1}{2}$)
  \item Alkaline earths: singlets and triplets
  \item Each multiplicity $= 2J+1$
\end{itemize}
\end{column}
\begin{column}{0.48\textwidth}
\textbf{Class~3: Anomalous Zeeman component counts}

\medskip
\begin{tabular}{p{3.4cm}cc}
\toprule
Na transition & Observed & Classical \\
\midrule
$^2P_{3/2}\to{}^2S_{1/2}$ (D$_2$) & \textbf{6} & 3 \\
$^2P_{1/2}\to{}^2S_{1/2}$ (D$_1$) & \textbf{4} & 3 \\
\bottomrule
\end{tabular}

\medskip
\small D$_2$: upper state $J=3/2$ (4 sublevels),
lower $J=1/2$ (2 sublevels),
$\Delta m_J \in \{-1,0,+1\}$ $\Rightarrow$ \textbf{6 lines}.

\bigskip
\begin{alertblock}{The Anomaly}
6 and 4 lines (not~3) are
\textbf{impossible} with integer $J$.\\
The four-parameter model gives the
correct count deductively.
\end{alertblock}
\end{column}
\end{columns}

\end{frame}

%----------------------------------------------------------------------
\begin{frame}[shrink=50]{Goodness of Fit: Zero Residuals}

\begin{center}
\renewcommand{\arraystretch}{1.2}
\begin{tabular}{cccrrc}
\toprule
Subshell & $l$ & Formula $2(2l+1)$ &
Predicted & Observed & Residual \\
\midrule
$1s$ & 0 & $2(1)$ &  2 &  2 & 0 \\
$2s$ & 0 & $2(1)$ &  2 &  2 & 0 \\
$2p$ & 1 & $2(3)$ &  6 &  6 & 0 \\
$3s$ & 0 & $2(1)$ &  2 &  2 & 0 \\
$3p$ & 1 & $2(3)$ &  6 &  6 & 0 \\
$3d$ & 2 & $2(5)$ & 10 & 10 & 0 \\
$4s$ & 0 & $2(1)$ &  2 &  2 & 0 \\
$4p$ & 1 & $2(3)$ &  6 &  6 & 0 \\
$4d$ & 2 & $2(5)$ & 10 & 10 & 0 \\
$4f$ & 3 & $2(7)$ & 14 & 14 & 0 \\
\midrule
\multicolumn{4}{l}{RSS} & & $\mathbf{0}$ \\
\bottomrule
\end{tabular}
\end{center}

\[
  \chi^2 = \sum_{\text{subshells}}
  \frac{(\text{Observed} - \text{Predicted})^2}{\text{Predicted}} = 0
\]

\begin{block}{Compare with Einstein}
Einstein's five-point sodium regression: $R^2 \approx 0.99999$,
residuals $\leq 2$\,mV.  Pauli: $\chi^2 = 0$.  The counts are
\emph{derived}, not \emph{fitted}.
\end{block}

\end{frame}

%----------------------------------------------------------------------
\begin{frame}[shrink=40]{Parameter Identification}

\textbf{Which observation pins down which parameter?}

\bigskip
\begin{center}
\renewcommand{\arraystretch}{1.2}
\begin{tabular}{p{3.2cm}p{4.5cm}p{5.5cm}}
\toprule
\textbf{Parameter} & \textbf{What it determines} & \textbf{Identifying observation} \\
\midrule
$n\in\{1,2,3,\ldots\}$ & Shell index; $n$ subshell types &
  Chemical periodicity; discrete X-ray absorption edges \\[4pt]
$l\in\{0,\ldots,n-1\}$ & Subshell; $2l+1$ magnetic levels &
  Fine structure of X-ray binding energies; Sommerfeld \\[4pt]
$m_l\in\{-l,\ldots,+l\}$ & Normal Zeeman splitting into $2l+1$ lines &
  Normal Zeeman effect; integer-spaced triplets \\[4pt]
$m_s\in\{\pm\tfrac{1}{2}\}$ & Factor-of-2 per sublevel; doublet spectra &
  Stoner's factor of~2; alkali doublets; Landé $g_s=2$ \\[4pt]
Exclusion $n_s\leq 1$ & Total capacity $= 2n^2$ &
  Noble-gas inertness; absence of all electrons collapsing to $1s$ \\
\bottomrule
\end{tabular}
\end{center}

\bigskip
\begin{alertblock}{Key asymmetry}
Parameters $n$, $l$, $m_l$ were already known.  Only~$m_s$ is new---and it
is identified \emph{indirectly}, through its consequences (the factor of~2),
not through a direct measurement.
\end{alertblock}

\end{frame}

%----------------------------------------------------------------------
\begin{frame}[shrink=20]{A Multinomial Likelihood Framing}

\textbf{Treat an electron drawn uniformly from shell $n$ as an observation.}

\medskip
Under Pauli's model, the probability it lands in subshell~$l$ is
\[
  p^*_l = \frac{2(2l+1)}{2n^2} = \frac{2l+1}{n^2},
  \quad l = 0,1,\ldots,n-1.
\]

\medskip
\begin{columns}[T]
\begin{column}{0.52\textwidth}
\textbf{$\chi^2$ test vs.\ flat null} ($p^{\text{null}}_l = 1/n$):

\smallskip
\begin{tabular}{cccr}
\toprule
Shell & Subshell & Pred.\ vs.\ null & $\chi^2$ \\
\midrule
$n=2$ & $s$ ($l=0$) & 2 vs.\ 4 & 1.00 \\
$n=2$ & $p$ ($l=1$) & 6 vs.\ 4 & 1.00 \\
\midrule
\multicolumn{2}{r}{$n=2$ total} & & 2.00 \\
\midrule
$n=3$ & $s$ ($l=0$) & 2 vs.\ 6 & 2.67 \\
$n=3$ & $p$ ($l=1$) & 6 vs.\ 6 & 0.00 \\
$n=3$ & $d$ ($l=2$) & 10 vs.\ 6 & 2.67 \\
\midrule
\multicolumn{2}{r}{$n=3$ total} & & 5.33 \\
\bottomrule
\end{tabular}
\end{column}
\begin{column}{0.45\textwidth}
\begin{block}{Parsimony}
The null model of equal subshell
occupancy uses $n-1$ free parameters
per shell.

\smallskip
Pauli's model uses \textbf{zero}.

\smallskip
AIC favours Pauli by $2(n-1)$ units
\emph{plus} the likelihood gain from
exact fit.
\end{block}

\medskip
\begin{alertblock}{Generalisation}
Pauli's formula extends analytically
to any $n$---shells 5, 6, 7 of xenon
and the actinides---without refitting.
A purely empirical model fitted to
$n=1,\ldots,4$ has no such basis.
\end{alertblock}
\end{column}
\end{columns}

\end{frame}

%----------------------------------------------------------------------
\begin{frame}[shrink=15]{Subsequent Confirmation: Spin (Uhlenbeck--Goudsmit, 1925)}

\begin{columns}[T]
\begin{column}{0.55\textwidth}
In September 1925---eight months after Pauli's paper---Uhlenbeck and
Goudsmit proposed that the electron has an \textbf{intrinsic angular
momentum (spin)} $s = \hbar/2$ with projection $m_s = \pm\tfrac{1}{2}$
and an anomalous gyromagnetic ratio $g_s \approx 2$.

\bigskip
\textbf{What this gave Pauli's model:}
\begin{itemize}
  \item A \emph{physical object} for the abstract $m_s$
  \item An explanation of the doublet structure of alkali spectra
    via spin--orbit coupling $\mathbf{L}\cdot\mathbf{S}$
  \item A derivation of the Landé g-factor formula with $g_s = 2$
\end{itemize}

\bigskip
Pauli was initially \textbf{sceptical}: a classically spinning electron
the size of a point particle would need super-luminal surface velocity.
He came to accept spin only after Dirac (1928).
\end{column}
\begin{column}{0.42\textwidth}
\begin{block}{Timeline}
\begin{tabular}{l p{4cm}}
Jan.\ 1925 & Pauli: exclusion principle; $m_s$ as formal device \\
Sep.\ 1925 & Uhlenbeck--Goudsmit: $m_s$ = intrinsic spin \\
1926 & Fermi--Dirac statistics \\
1928 & Dirac: spin from relativity \\
1940 & Pauli: spin--statistics theorem \\
1945 & Pauli Nobel Prize \\
\end{tabular}
\end{block}
\end{column}
\end{columns}

\end{frame}

%----------------------------------------------------------------------
\begin{frame}[shrink=20]{Fermi--Dirac Statistics and the Exclusion Principle}

\textbf{In 1926, Fermi and Dirac independently derived the quantum statistics
of particles obeying the exclusion principle:}

\[
  \bar{n}(\varepsilon) = \frac{1}{e^{(\varepsilon-\mu)/k_BT} + 1}
\]

\begin{columns}[T]
\begin{column}{0.52\textwidth}
The \textbf{$+1$} in the denominator is the signature of the Exclusion Principle:
$0 < \bar n < 1$ for all~$\varepsilon$.

\medskip
At $T=0$: step function.  All states below Fermi energy $E_F$ are full;
all above are empty.

\bigskip
\textbf{Immediate applications:}
\begin{itemize}
  \item Fowler (1926): \emph{white dwarf stability}---electron degeneracy
    pressure resists gravitational collapse
  \item Sommerfeld (1927): \emph{metallic heat capacity}---at room temperature
    only electrons within $k_BT$ of $E_F$ can be excited
\end{itemize}
\end{column}
\begin{column}{0.45\textwidth}
\begin{alertblock}{Bose--Einstein vs.\ Fermi--Dirac}
\[
  \bar n_{\text{BE}} = \frac{1}{e^{(\varepsilon-\mu)/k_BT} - 1}
  \quad (\text{bosons})
\]
\[
  \bar n_{\text{FD}} = \frac{1}{e^{(\varepsilon-\mu)/k_BT} + 1}
  \quad (\text{fermions})
\]
The $\pm 1$ encodes the exclusion or
accumulation physics.
\end{alertblock}

\bigskip
\begin{block}{Transfer to Astrophysics}
The same counting rule that governs
atoms explains why white dwarf stars
do not collapse.
\end{block}
\end{column}
\end{columns}

\end{frame}

%----------------------------------------------------------------------
\begin{frame}[shrink=15]{Dirac (1928) and the Spin--Statistics Theorem (1940)}

\begin{columns}[T]
\begin{column}{0.52\textwidth}
\textbf{Dirac (1928):} relativistic quantum equation for the electron.
\medskip
\begin{itemize}
  \item Spin-$\tfrac{1}{2}$ emerges \emph{automatically} from combining
    QM with special relativity---not added by hand
  \item $g_s = 2$ exactly (including corrections from QED: $g_s = 2.00232\ldots$)
  \item Retroactively \emph{derives} Pauli's fourth quantum number
    from first principles
\end{itemize}

\bigskip
\textbf{Pauli's spin--statistics theorem (1940):}\\
In any relativistic QFT with locality,
causality, and positive energy:
\begin{itemize}
  \item Half-integer spin particles \textbf{must} obey the Exclusion Principle
  \item Integer spin particles must obey Bose--Einstein statistics
\end{itemize}
\end{column}
\begin{column}{0.45\textwidth}
\begin{block}{From Postulate to Theorem}
\begin{tabular}{p{1.5cm}p{4cm}}
\toprule
1925 & Exclusion Principle: a \emph{postulate} with no derivation \\
1928 & Spin from Dirac equation: fourth quantum number derived \\
1940 & Spin--statistics theorem: Exclusion Principle is a \emph{theorem} \\
\bottomrule
\end{tabular}
\end{block}

\bigskip
\begin{alertblock}{Nobel Prize 1945}
Awarded to Pauli ``for the discovery of
the Exclusion Principle, also called the
Pauli Principle.''
\end{alertblock}
\end{column}
\end{columns}

\end{frame}

%----------------------------------------------------------------------
\begin{frame}[shrink=30]{The ML-to-Physics Dictionary}

\begin{center}
\footnotesize
\renewcommand{\arraystretch}{1.10}
\begin{tabular}{p{5.5cm}p{4.8cm}p{3.5cm}}
\toprule
\textbf{Scientific step} & \textbf{ML name} & \textbf{Who / when} \\
\midrule
Capacity $= 2n^2$, not $n^2$; anomalous Zeeman lines &
  Anomalous training-set residuals & Bohr--Sommerfeld model by 1922 \\
Stoner: factor of~2 per $m_l$ sublevel &
  Feature-importance / residual diagnosis & Stoner (1924) \\
Propose $m_s\in\{+\tfrac{1}{2},-\tfrac{1}{2}\}$;
  exclusion $n_s\leq 1$ &
  Inductive bias + hypothesis class & Pauli (1925) \\
State space: $(n,l,m_l,m_s)$ lattice &
  Feature engineering (add 4th feature) & Pauli (1925) \\
Count states per subshell: $2(2l+1)$ &
  Forward pass of model & Pauli (1925) \\
$\chi^2 = 0$: all 10 capacity observations fit &
  Test accuracy / Confirmation & Spectroscopy 1913--1924 \\
Transferable to molecules, nuclei, stars &
  Zero-shot transfer learning & Fowler (1926); Sommerfeld (1927) \\
Spin--statistics theorem &
  Theoretical derivation of architecture & Pauli (1940) \\
\bottomrule
\end{tabular}
\end{center}

\end{frame}

%----------------------------------------------------------------------
\begin{frame}[shrink=25]{The Hypothesis Class: Where Pauli Differs from a Naive ML System}

\begin{columns}[T]
\begin{column}{0.52\textwidth}
\textbf{Pauli's hypothesis class:}
\[
  \mathcal{H} = \bigl\{\text{binary occupation on }
  (n,l,m_l,m_s)\text{ lattice}\bigr\}
\]
with $n_s \leq 1$.

\medskip
\begin{itemize}
  \item \textbf{VC dimension $= 1$}: only one model in the class
  \item Single observation that confirms $2(2l+1)$ suffices theoretically
  \item Zero adjustable parameters
  \item Derived from two postulates, not from data
\end{itemize}
\end{column}
\begin{column}{0.45\textwidth}
\begin{alertblock}{What a Naive ML System Misses}
\begin{enumerate}
  \item Cannot infer from $(n,l) \to \text{count}$ data alone that
    the explanation is a \emph{two-valued} quantum number
  \item Cannot connect capacity data to Zeeman multiplicities without
    the shared latent variable (quantum number structure)
  \item Cannot predict Fermi--Dirac statistics or white dwarf stability
    from the atomic table
  \item Cannot generalise to $n=5,6,7$ without Pauli's analytic formula
\end{enumerate}
\end{alertblock}
\end{column}
\end{columns}

\bigskip
\begin{block}{The Compression Pauli Achieved}
$m_s + n_s\leq 1$ collapsed the hypothesis space from all occupation rules
on an arbitrarily rich state space to a \textbf{zero-parameter counting function}---
the most extreme inductive bias possible.
\end{block}

\end{frame}

%----------------------------------------------------------------------
\begin{frame}[shrink=40]{Pauli vs.\ a Deep Neural Network}

\begin{center}
\small
\renewcommand{\arraystretch}{1.2}
\begin{tabular}{p{3.0cm}p{2.5cm}p{3.5cm}p{4.5cm}}
\toprule
Epoch & Data available & Pauli model & Neural-network outcome \\
\midrule
\textbf{Bohr--Sommerfeld} (1922) &
  Qualitative periodicity &
  Diagnoses: 3 QN give $n^2$ not $2n^2$ &
  Sees integers; cannot diagnose missing feature \\[3pt]
\textbf{Stoner} (1924) &
  Subshell counts; 10 integers &
  Identifies binary 4th QN; zero parameters &
  Could fit a regression on 10 points, but has no basis for $2(2l+1)$ form \\[3pt]
\textbf{Pauli} (1925) &
  As above + spectral doublets &
  Predicts all observations; zero free parameters &
  Memorises; cannot extrapolate to $n=5$ or Zeeman counts \\[3pt]
\textbf{Modern} (large corpus) &
  All elements + spectroscopy &
  Same formula; no retraining needed &
  Could learn $2(2l+1)$ numerically; cannot identify it as exclusion principle \\
\bottomrule
\end{tabular}
\end{center}

\bigskip
\begin{alertblock}{The Core Lesson}
Pauli's ``training set'' was \textbf{10 integer observations}.
His model has \textbf{zero free parameters}.
A neural network cannot make use of 10 observations---it needs theory.
\end{alertblock}

\end{frame}

%----------------------------------------------------------------------
\begin{frame}[shrink=15]{Zero-Shot Transfer: From Atoms to Stars}

\begin{columns}[T]
\begin{column}{0.52\textwidth}
\textbf{The exclusion principle transfers across all fermion systems:}

\medskip
\begin{tabular}{p{2.2cm}p{5.5cm}}
\toprule
Domain & How Exclusion Principle applies \\
\midrule
Atoms & Shell capacities $2n^2$; chemical periodicity \\
Molecules & Hückel's $4n+2$ rule for aromatic stability \\
Solids & Fermi surface; low heat capacity of metals \\
White dwarfs & Electron degeneracy pressure prevents collapse \\
Neutron stars & Neutron degeneracy pressure prevents collapse \\
Nuclear physics & Nuclear shell model; magic numbers \\
\bottomrule
\end{tabular}
\end{column}
\begin{column}{0.45\textwidth}
\begin{alertblock}{ML Language}
\textbf{Pre-training:} Exclusion Principle
derived from atomic data\\[4pt]
\textbf{Zero-shot transfer:} same principle
applied to electrons in metals, nuclear
protons/neutrons, stellar interiors\\[4pt]
\textbf{Fine-tuning needed:} only the
energy spectrum $\varepsilon_i$ relevant to
each domain---the occupation rule
transfers unchanged
\end{alertblock}

\bigskip
\begin{block}{Compare with Einstein}
Einstein's slope $h/e$: Na to Li,
0.2\% agreement.\\
Pauli's exclusion rule: atoms to stars,
\emph{exact}.
\end{block}
\end{column}
\end{columns}

\end{frame}

%----------------------------------------------------------------------
\begin{frame}[shrink=20]{Pauli's Epistemic Situation}

\begin{columns}[T]
\begin{column}{0.52\textwidth}
Pauli proposed the exclusion principle and the fourth quantum number in
January 1925.  \textbf{Then wrote:}

\bigskip
\begin{quote}
``Already in my original paper I stressed the circumstance that I was
\textbf{unable to give a logical reason} for the exclusion principle
or to derive it from more general assumptions\ldots\ To me, the
suggestion of Stoner appeared to be an important step forward, since
for the first time \textbf{the number~2} had to be attributed to the
electron itself---to its `two-valuedness.'\,''
\end{quote}
\hfill\small{--- Pauli, Nobel Lecture (1946)}
\end{column}
\begin{column}{0.45\textwidth}
\begin{block}{Lesson for Statistics}
A parsimonious formal model can unify a body of empirical observations
\textbf{before} the physical mechanism is understood.

\bigskip
Pauli in 1925 did not know what $m_s$
was.  In this he resembles:
\begin{itemize}
  \item Planck (1900): introduced $h\nu$
    without knowing what a ``quantum'' was
  \item Einstein (1905): proposed light
    quanta that conflicted with wave theory
  \item Millikan (1916): confirmed $h/e$
    while rejecting the light-quantum model
\end{itemize}
\end{block}
\end{column}
\end{columns}

\end{frame}

%----------------------------------------------------------------------
\begin{frame}[shrink=20]{Comparison with the Photoelectric Episode}

\begin{center}
\small
\renewcommand{\arraystretch}{1.15}
\begin{tabular}{p{3.3cm}p{5.0cm}p{5.0cm}}
\toprule
\textbf{Dimension} & \textbf{Einstein (photoelectric)} & \textbf{Pauli (exclusion principle)} \\
\midrule
Anomaly & Stopping potential independent of intensity &
  Shell capacity $2n^2$ instead of $n^2$ \\
Key precursor & Lenard (1902): $V_0 \propto \nu$ qualitatively &
  Stoner (1924): factor-of-2 per sublevel \\
Postulate & $E = h\nu$ (light quantum) &
  $m_s\in\{\pm\tfrac{1}{2}\}$ + $n_s\leq 1$ \\
Model type & Continuous regression & Discrete counting \\
Free parameters & 2 ($\beta=h/e$; $\alpha=-W/e+v_c$) & \textbf{0} \\
Fit statistic & $R^2 \approx 0.99999$ (5 points) & $\chi^2 = 0$ (10 integers) \\
Transfer & Na to Li, 0.2\% & Atoms to stars, exact \\
Interpretation delayed & Light-quantum debated until 1923 &
  $m_s$ unidentified until Uhlenbeck--Goudsmit (1925) \\
Grounding achieved & QED (1940s) &
  Dirac (1928); Pauli spin--statistics (1940) \\
Nobel Prize & 1921 (Einstein) & 1945 (Pauli) \\
\bottomrule
\end{tabular}
\end{center}

\end{frame}

%----------------------------------------------------------------------
\begin{frame}[shrink=55]{The Hypothetico-Deductive Loop}

\begin{center}
\small
\renewcommand{\arraystretch}{1.25}
\begin{tabular}{p{4.2cm}p{4.2cm}p{4.2cm}}
\toprule
\textbf{Step} & \textbf{Science} & \textbf{ML} \\
\midrule
Diagnose $2n^2 \neq n^2$ & Empirical anomaly & Training residuals \\
Stoner: factor-of-2 per sublevel & Key feature identified & Feature importance \\
Propose $m_s\in\{\pm\tfrac{1}{2}\}$; $n_s\leq 1$ & Theoretical hypothesis & Model architecture \\
State space: $(n,l,m_l,m_s)$ & Deduction & Hypothesis class \\
Count $2(2l+1)$ per subshell & Prediction & Forward pass \\
$\chi^2 = 0$: all 10 counts exact & Confirmation & Test accuracy \\
Fermi--Dirac, white dwarfs, nuclei & Universality & Zero-shot transfer \\
\bottomrule
\end{tabular}
\end{center}

\bigskip
\begin{block}{Modern AI for Science}
An AI Feynman-type system given the 10-integer table could find the formula
$2(2l+1)$ by symbolic search.  What it cannot do:
\begin{itemize}
  \item Explain \emph{why}---the exclusion principle as a law of nature
  \item Connect the formula to spectral doublets and Zeeman multiplicities
  \item Transfer the rule to white dwarfs without re-learning
\end{itemize}
\textbf{Theory construction---not number fitting---is Pauli's irreducible contribution.}
\end{block}

\end{frame}

%----------------------------------------------------------------------
\begin{frame}[shrink=35]{Summary and Conclusions}

\begin{block}{The Thesis, Revisited}
Pauli in 1925 was doing machine learning---but his decisive step was the
one no ML system has yet automated: \textbf{specifying the hypothesis
class} (what a state space looks like and how many electrons it can hold).
\end{block}

\bigskip
\textbf{Three principles illustrated:}
\begin{enumerate}
  \item \textbf{Extreme parsimony.}  One postulate---$n_s\leq 1$---organises
    the entire periodic table, all atomic spectral multiplicities, Fermi--Dirac
    statistics, and stellar stability.  Zero free parameters.
  \item \textbf{Theory-first identification.}  The fourth quantum number was
    diagnosed from a qualitative discrepancy (factor of~2), not from a new
    precision measurement.  Confirmation followed theory.
  \item \textbf{Uninterpreted success.}  Pauli proposed $m_s$ without knowing
    its physical meaning.  Planck, Einstein, and Pauli each introduced a
    formal device that organised data parsimoniously, long before a
    self-consistent interpretation was available.
\end{enumerate}

\bigskip
\begin{alertblock}{The Limit of Pure AI}
A data-driven system can recover $2(2l+1)$ numerically but cannot know it
is the \emph{exclusion principle}.
\textbf{Physical meaning requires physical theory.}
\end{alertblock}

\end{frame}

%----------------------------------------------------------------------
\begin{frame}[allowframebreaks]{References}

\small
\begin{itemize}
  \item Pauli, W.\ (1925). ``\"{U}ber den Zusammenhang des Abschlusses der
    Elektronengruppen im Atom mit der Komplexstruktur der Spektren.''
    \textit{Z.\ Phys.}, 31, 765--783.
  \item Stoner, E.\,C.\ (1924). ``The Distribution of Electrons among
    Atomic Levels.''
    \textit{Phil.\ Mag.}, 48(286), 719--736.
  \item Uhlenbeck, G.\,E.\ \& Goudsmit, S.\ (1925).
    ``Ersetzung der Hypothese vom unmechanischen Zwang\ldots''
    \textit{Naturwissenschaften}, 13, 953--954.
  \item Land\'{e}, A.\ (1921). ``\"{U}ber den anomalen Zeemaneffekt.''
    \textit{Z.\ Phys.}, 5, 231--241.
  \item Fermi, E.\ (1926). ``Sulla quantizzazione del gas perfetto
    monoatomico.'' \textit{Rend.\ Lincei}, 3, 145--149.
  \item Dirac, P.\,A.\,M.\ (1926). ``On the Theory of Quantum Mechanics.''
    \textit{Proc.\ R.\ Soc.\ A}, 112, 661--677.
  \item Dirac, P.\,A.\,M.\ (1928). ``The Quantum Theory of the Electron.''
    \textit{Proc.\ R.\ Soc.\ A}, 117, 610--624.
  \item Pauli, W.\ (1940). ``The Connection Between Spin and Statistics.''
    \textit{Phys.\ Rev.}, 58, 716--722.
  \item Pauli, W.\ (1946). ``Exclusion Principle and Quantum Mechanics.''
    Nobel Lecture, December~13, 1946.
    In: \textit{Nobel Lectures, Physics 1942--1962}, Elsevier, pp.~27--43.
  \item Fowler, R.\,H.\ (1926). ``On Dense Matter.''
    \textit{Mon.\ Not.\ R.\ Astron.\ Soc.}, 87, 114--122.
  \item Sommerfeld, A.\ (1927). ``Zur Elektronentheorie der Metalle\ldots''
    \textit{Z.\ Phys.}, 47, 1--32.
  \item Massimi, M.\ (2005). \textit{Pauli's Exclusion Principle}.
    Cambridge University Press.
  \item Udrescu, S.\,M.\ \& Tegmark, M.\ (2020). ``AI Feynman.''
    \textit{Sci.\ Adv.}, 6, eaay2631.
\end{itemize}

\end{frame}

%----------------------------------------------------------------------
\end{document}
